
sera obligé de travailler avec l'anneau de Chow, puisque c'est sur ce dernier que se trouve défini naturellement l'homomorphisme \(\operatorname{c\ell}\) comme homomorphisme à valeurs dans \(H^{2*}(-,\mathbb Z_\ell(*))\), comparer (XIV 4.1). Malheureusement, on n'a pas développé pour le moment une théorie de la spécialisation pour l'anneau de Chow (cf. 7.14.), qui permettrait d'écrire des carrés commutatifs analogues à (7.13.9) et (7.13.10) en termes des anneaux de Chow et de la cohomologie à coefficients \(\mathbb Z_\ell(*)\). Tout ce qu'on peut dire lorsque \(f:X\to S\) est lisse et propre, c'est qu'on a un carré commutatif
\begin{equation}\tag{7.13.12}
\begin{tikzcd}[column sep=large,row sep=large]
Z^i(X_t) \arrow[r,"\operatorname{c\ell}"] \arrow[d,"\sigma"'] &
H^{2i}(X_t,\mathbb Z_\ell(i))\\
Z^i(X_s) \arrow[r,"\operatorname{c\ell}"] &
H^{2i}(X_s,\mathbb Z_\ell(i)) \arrow[u]
\end{tikzcd}
\;,
\end{equation}
analogue à (7.13.9), et un carré correspondant analogue à (7.13.10) qu'on laisse au lecteur le soin d'écrire, où \(Z^i\) désigne le groupe des cycles algébriques (modulo rien du tout) de codimension \(i\), où les flèches horizontales sont les homomorphismes «~classe de cohomologie associée~» définis dans (SGA 5 IV), et où \(\sigma\) désigne l'homomorphisme de spécialisation des cycles algébriques, associant à un cycle \(Z_t\) le cycle \(Z_s\), où \(Z\) est l'adhérence schématique de \(Z_t\) dans \(X\) (cf. 7.14). Cette commutativité résulte du fait suivant, établi dans (SGA 5 IV); si \(Z\) est un sous-schéma fermé de \(X\) de codimension \(\geq i\) sur chaque fibre, on peut lui associer une classe de cohomologie \(\operatorname{c\ell}(Z)\in H^{2i}(X,\mathbb Z_\ell(i))\), dont la formation commute à tout changement de base, et en particulier telle que sa restriction à \(X_t\) (resp. \(X_s\)) soit égale à \(\operatorname{c\ell}(Z_t)\) (resp. \(\operatorname{c\ell}(Z_s)\)). Bien entendu, lorsqu'on sait que l'homomorphisme \(\sigma\) dans (7.13.12) transforme cycles rationnellement équivalents à zéro en cycles rationnellement équivalents à zéro (\(X\) étant projectif sur \(S\)),
